% !PRINTED 750
\weberpage{750}

Nun ergiebt sich sehr einfach, dass keine Zahl $g$, die mit einer der Zahlen $a, b$ zusammenf\"allt, oder auch, wenn $a$ und $b$ verschieden sind, zwischen ihnen liegt, zu der Reihe $\Omega$ geh\"oren kann.

% [sic] -- the source prints singular ``kommt'' with the pair $\alpha_\nu,\beta_\nu$; the printed German is preserved.
Denn die Zahl $g$ liegt im Inneren eines jeden der Intervalle $\delta_\nu$. Nehmen wir an, es komme $g$ in $\Omega$ vor, und gehen nach 3. mit $\nu$ so weit, dass $\alpha_\nu, \beta_\nu$ in $\Omega$ sp\"ater als $g$ kommt, so kann $g$ nach 2. nicht im Intervall $\delta_\nu$ liegen, und damit ist die Unm\"oglichkeit unserer Annahme erwiesen.

\gesperrt{Daraus folgt also, dass die Gesammtheit der Zahlen eines Intervalles $\alpha, \beta$ keine abz\"ahlbare Menge bildet.}

\begin{center}
\rule{4cm}{0.4pt}
\end{center}

Diese Thatsache l\"asst sich noch auf einem anderen Wege beweisen, der in gewisser Beziehung noch einfacher ist und den wir mit wenig Worten darlegen wollen. Wir beschr\"anken die Allgemeinheit nicht, wenn wir das Intervall von 0 bis 1 zu Grunde legen. Alle Zahlen dieses Intervalles denken wir uns durch unendliche Decimalbr\"uche dargestellt. Darunter sind auch die endlichen Decimalbr\"uche enthalten, wenn wir alle Ziffern, von einer gewissen an, gleich Null setzen. Um die Darstellung durch Decimalbr\"uche zu einer eindeutigen zu machen, mag noch festgesetzt sein, dass f\"ur einen endlichen Decimalbruch \gesperrt{immer} diese Darstellung gew\"ahlt, also z.\,B. nicht 0,4999\dots\ f\"ur 0,5000\dots\ gesetzt werden soll.

Wir wollen nun annehmen, diese Decimalbr\"uche bilden eine abz\"ahlbare Menge; sie lassen sich also in eine z\"ahlbare Reihe anordnen, die wir so darstellen:
\weq{$\Omega$}{
\begin{aligned}
\omega_1 &= 0,\ \alpha_1^{(1)}\, \alpha_2^{(1)}\, \alpha_3^{(1)}\, \dots\\
\omega_2 &= 0,\ \alpha_1^{(2)}\, \alpha_2^{(2)}\, \alpha_3^{(2)}\, \dots\\
\omega_3 &= 0,\ \alpha_1^{(3)}\, \alpha_2^{(3)}\, \alpha_3^{(3)}\, \dots\\
&\quad\cdots\cdots\cdots\cdots\cdots
\end{aligned}
}
worin die $\alpha_\mu^{(\nu)}$ Ziffern des dekadischen Systems bedeuten.

Es ist nun aber sehr leicht, einen Decimalbruch (oder auch beliebig viele) nachzuweisen, die in der Reihe $\Omega$ nicht enthalten sind. Wir brauchen nur
\[
\eta = 0,\ \beta_1\ \beta_2\ \beta_3\ \dots
\]
% [sic] -- the printed condition permits a recurring-9 expansion equivalent to an earlier terminating decimal; one repair is to choose $\beta_\nu\in\{1,2\}$ unequal to $\alpha_\nu^{(\nu)}$.
zu bilden, wobei die $\beta_\nu$ Ziffern des dekadischen Systems sind, die der einen Bedingung gen\"ugen, dass $\beta_\nu$ f\"ur jedes $\nu$ von $\alpha_\nu^{(\nu)}$ ver-
%CONT>>
